% SEGMENT OLP-0111-S01
% Part: first-order-logic
% Chapter: tableaux
% Section: soundness-identity

\documentclass[../../../include/open-logic-section]{subfiles}

\begin{document}

\olfileid{fol}{tab}{sid}

\olsection{\usetoken{S}{identity} கொண்ட சரித்தன்மை}

\begin{prop}
ஒன்றாதலுக்கான விதிகளைக் கொண்ட !!^{tableau}s சரித்தன்மையுடையவை: எந்த மூடிய
!!{tableau}வும் நிறைவுறத்தக்கதல்ல.
\end{prop}

% SEGMENT OLP-0111-S02
\begin{proof}
முன்புபோலப் பின்வருவதை மட்டும் காட்ட வேண்டும்: !!a{tableau} ஒன்றுக்கு
நிறைவுறத்தக்க கிளை இருக்கும்போது, அந்தக் கிளையில்~$\eq$ க்கான விதிகளுள்
ஒன்றைப் பயன்படுத்துவதால் கிடைக்கும் கிளையும் நிறைவுறத்தக்கது. கிளையிலுள்ள
குறியிட்ட !!{formula}s கணத்தை~$\Gamma$ எனவும், $\Gamma$ வை நிறைவுறுத்தும்
!!a{structure} ஒன்றை~$\Struct{M}$ எனவும் கொள்க.

$\eq$ ஐப் பயன்படுத்தி, அதாவது குறியிட்ட
!!{formula}~$\sFmla{\True}{\eq[t][t]}$ ஐச் சேர்த்து, கிளை
விரிவாக்கப்படுகிறது என்க. தெளிவாகவே $\Sat{M}{\eq[t][t]}$; ஆகவே
$\Struct{M}$ ஆனது $\Gamma \cup \{\sFmla{\True}{\eq[t][t]}\}$ ஐயும்
நிறைவுறுத்துகிறது.

% SEGMENT OLP-0111-S03
$\TRule{\True}{\eq}$ ஐப் பயன்படுத்திக் கிளை விரிவாக்கப்பட்டால், குறியிட்ட
% SOURCE CORRECTION TA-TAB-006: the T-identity conclusion has the True sign.
!!{formula}~$\sFmla{\True}{!A(t_2)}$ ஐச் சேர்க்கிறோம்; மேலும் $\Gamma$ ஆனது
$\sFmla{\True}{\eq[t_1][t_2]}$, $\sFmla{\True}{!A(t_1)}$ ஆகிய இரண்டையும்
கொண்டுள்ளது. எனவே $\Sat{M}{\eq[t_1][t_2]}$ மற்றும்
$\Sat{M}{!A(t_1)}$. $s(x) = \Value{t_1}{M}$ ஆகுமாறு ஒரு மாறி
ஒதுக்கீட்டை~$s$ என்க.
\oliflabeldef{syn:ass:prop:sentence-sat-true}{\olref[syn][ass]{prop:sentence-sat-true}}{ஒரு
வாக்கியத்தின் நிறைவுறுத்தல் மாறி ஒதுக்கீட்டைச் சாராது என்ற முன்மொழிவு} படி,
$\Sat{M}{!A(t_1)}[s]$. $\varAssign{s}{s}{x}$ என்பதால்,
\oliflabeldef{syn:ext:prop:ext-formulas}{\olref[syn][ext]{prop:ext-formulas}}{சொல்
இடைநிறுத்தலுக்கும் மாறி ஒதுக்கீட்டுக்கும் இடையிலான முன்மொழிவு} படி,
$\Sat{M}{!A(x)}[s]$. $\Sat{M}{\eq[t_1][t_2]}$ என்பதால்,
$\Value{t_1}{M} = \Value{t_2}{M}$; எனவே $s(x) = \Value{t_2}{M}$.
\oliflabeldef{syn:ext:prop:ext-formulas}{\olref[syn][ext]{prop:ext-formulas}}{சொல்
இடைநிறுத்தலுக்கும் மாறி ஒதுக்கீட்டுக்கும் இடையிலான முன்மொழிவு} ஐ மீண்டும்
பயன்படுத்தினால், $\Sat{M}{!A(t_2)}[s]$ என்றும் கிடைக்கிறது.
\oliflabeldef{syn:ass:prop:sentence-sat-true}{\olref[syn][ass]{prop:sentence-sat-true}}{ஒரு
வாக்கியத்தின் நிறைவுறுத்தல் மாறி ஒதுக்கீட்டைச் சாராது என்ற முன்மொழிவு} படி,
$\Sat{M}{!A(t_2)}$. $\TRule{\False}{\eq}$ நேர்வும் இதேபோல்
கையாளப்படுகிறது.
\end{proof}

\end{document}
